(1) 纯气体初始化学势为 \(\mu_i^\circ(T,p)\)，混合后组分分压 \(p_i=x_ip\)，因此 \(\mu_i=\mu_i^\circ+RT\ln x_i\)。两态 Gibbs 自由能差为 \(\Delta G_{\rm mix}=n_ART\ln x_A+n_BRT\ln x_B\)。每个 \(x_i<1\)，故自由能严格下降，符合等温恒压下混合的自发趋势。

(2) 理想混合物分子能量无额外相互作用项，所以 \(\Delta H_{\rm mix}=0\)。由 \(\Delta G=\Delta H-T\Delta S\) 得 \(\Delta S_{\rm mix}=-R(n_A\ln x_A+n_B\ln x_B)>0\)。把 \(x_i\ln x_i\) 用 Gibbs 不等式或 \(0<x_i<1\) 的对数符号估计，亦能直接验证混合熵正定。

(3) 理想半透膜使两组分在各自化学势差下可逆分离，最低外部功等于逆向 Gibbs 自由能差，即 \(W_{\min}=-\Delta G_{\rm mix}=T\Delta S_{\rm mix}\)。可逆路径可逼近此下界，任何具有扩散、摩擦或不可逆混合的分离装置都需要不少于该值的实际外功。热力学势的使用依赖固定温度和压强，不可替换为随意的内能差。

(4) 若 \(n_B\to0\)，\(n_B\ln x_B\to0\)，并且 A 的摩尔分数趋一，混合熵趋零。若 \(n_A=n_B=n/2\)，有 \(\Delta S_{\rm mix}=nR\ln2\)。当 A、B 实际不可区分，是相同理想气体时移走隔板不产生新的可观测组分，正确熵变为零；将人为贴标签后使用双组分公式会造成 Gibbs 悖论。
